\section{The decade table}
\label{sec:table}

Our computations use the finite graph of~\cite{dbts} at the granularity of decades rather than
danger intervals, extended by the information needed for the child graph.

\begin{definition}
  Let $k \ge 2$. The \emph{inputs} of decade $k$ are the entries $b^k + u$ with $u \in E$ and the
  starts $c \cdot b^k$ with $2 \le c \le b - 1$. The \emph{passage} of an input is the successor
  path from it up to and including its first term in decade $k + 1$ (the \emph{exit}
  $b^{k+1} + u'$), or up to the landmine $b^{k+1} - r$ where it dies. We record the exit or
  landmine, the leading digits $d$ of the branch points on the way, and the numbers of steps.
\end{definition}

A passage consists of at most $b - 1$ stretches of constant leading digit. Let $M$ be the least
common multiple of all block increments $D_f(x)$; in base 10,
$M = 2^5 \cdot 3^3 \cdot 5^3 \cdot 7 \cdot 11 \cdot 13 \cdot 23 \cdot 29 = 72\,108\,036\,000$.

\begin{lemma}
  \label{lem:reduction}
  Let $k, k' \ge 3$ with $b^k \equiv b^{k'} \pmod M$. Then every input has the same exit or landmine
  in decade $k$ as in decade $k'$, and passes branch points with the same leading digits in the
  same order. The numbers of steps of its passages differ by $b\,(b^k - b^{k'}) \sum 1/D$, the sum
  running over the stretches of the passage, with $D$ the block increment of each stretch.
\end{lemma}

\begin{proof}
  Write $P = b^k$ and a term of stretch $f$ of decade $k$ as $n = fP + o$ with $0 \le o < P$, and
  let $t = P - o$ be its distance to the end of the stretch.

  \emph{The rule near a boundary.} As $b \mid P$, the last digit of $n$ is $x = (-t) \bmod b$. A
  child $n + xb + e$ is less than $n + b^2 \le n + P$, so its leading digit is $f$ if $xb + e < t$,
  and $f + 1$ otherwise, or $1$ in decade $k + 1$ if $f = b - 1$. Hence $n$ has the child with
  $e = f$ if and only if $xb + f < t$, the child with $e = f + 1$ (or $e = 1$ if $f = b - 1$) if
  and only if $xb + e \ge t$, and no other child. Which children exist, and where they lie relative
  to the boundary, therefore depends only on $f$ and $t$. In particular, a term with $t > b^2$ has
  the single child $n + xb + f$, and branch points and landmines have $t \le b^2$.

  \emph{The middle of a stretch.} A passage enters each stretch at an offset $o_0 < b^2$: the
  inputs have $o_0 = u \in E$ or $o_0 = 0$, and a child lies less than $b^2$ above its parent.
  So $t_0 = P - o_0 > b^3 - b^2 \ge b^2$. Let $x_0$ be the last digit on entry and $D = D_f(x_0)$.
  As long as the distance exceeds $b^2$, the terms are $n_0 + \sigma_j$, where $\sigma_j$ are the
  partial sums of the cycle of comma-numbers, with $\sigma_{j + b} = \sigma_j + D$. Let
  $q = \lfloor (t_0 - b^2 - 1) / D \rfloor \ge 0$. All terms up to $n_0 + qD$ have distance at
  least $t_0 - qD > b^2$, so the path reaches $n_0 + qD$ after $qb$ steps, with last digit $x_0$ and
  distance $t_0 - qD = b^2 + 1 + \bigl((t_0 - b^2 - 1) \bmod D\bigr)$. By the rule near a boundary,
  everything that follows in the stretch --- the remaining steps, a branch point or landmine on the
  way, and the offset at which the path enters the next stretch or decade --- depends only on $f$,
  $x_0$ and this distance, hence on $f$, $x_0$ and $t_0 \bmod D$.

  \emph{Comparing decades.} In decades $k$ and $k'$ an input enters its first stretch at the same
  offset, so the distances differ by $b^k - b^{k'}$, a multiple of $D$. By induction over the
  stretches, the passages agree stretch by stretch, and in each stretch $q$ differs by
  $(b^k - b^{k'}) / D$, that is, the number of steps by $b\,(b^k - b^{k'}) / D$.
\end{proof}

The second part of the proof is the argument of~\cite[Proposition~8]{dbts}, where the distance
$b^k + u$ is reduced modulo the sum of one cycle, $D_f(x)$ or a divisor of it. The lemma adds what we
need beyond their statement: an explicit range of $k$, the branch points, and the step counts. The
branch points could also be read off their points $(d, u, k)$: since the smaller child of the branch
point below $(d+1)\,b^k$ is $(d+1)\,b^k - 1$~\cite[Theorem~5.2]{comma}, a successor path passes it if
and only if $(d+1)\,b^k - 1$ is its last term below $(d+1)\,b^k$, i.e.\ $u = 1$. The step counts are
used only to compute term indices, not in the results of Section~\ref{sec:consequences}.

\begin{corollary}
  \label{cor:period}
  In base 10, the passages through decades $k$ and $k + 924$ agree, in the sense of
  Lemma~\ref{lem:reduction}, for all $k \ge 5$.
\end{corollary}

\begin{proof}
  For $k \ge 5$, $10^k$ is divisible by $2^5 \cdot 5^3$, and modulo the remaining factor
  $3^3 \cdot 7 \cdot 11 \cdot 13 \cdot 23 \cdot 29$ of $M$ it has period $924$, the lcm of the
  orders $3, 6, 2, 6, 22, 28$ of $10$ modulo $27, 7, 11, 13, 23, 29$.
\end{proof}

The threshold is sharp: some passages through decades $3$ and $4$ differ from those through
decades $927$ and $928$. The period $924$ is $L(10)$ of~\cite{dbts}.

\paragraph{Implementation.}
The proof of Lemma~\ref{lem:reduction} uses $b^k$ only through its residue modulo $M$, the fact that
it is a multiple of $b$, and $b^k - b^2 \ge b^2$. So it applies verbatim if $b^k$ is replaced by a
number $B \equiv b^k \pmod M$ with $B \ge 2b^2$, writing the numbers of a decade as $f B + o$, and
the implementation computes with such representatives; the step counts for $b^k$ follow from those
for $B$ as in the lemma, which makes term indices exact. Its choice of representatives repeats from
$k = 11$ on, so the table consists of the passages through the decades $11, \dots, 934$, one for each
of the $54$ inputs, $924 \cdot 54 = 49\,896$ in total; this is also the number of vertices counted by
the program of~\cite{dbtscode} for $b = 10$. By Corollary~\ref{cor:period}, the table also describes
the decades $5, \dots, 10$.

\paragraph{Validation.}
\begin{itemize}
  \item For bases $3, 4, 5, 6$ and $10$, the passages computed with the reduced $B$ agree with those
        computed with the real $b^k$ (exit, landmine, branch points and step counts) for the
        first $30$ decades of the periodic range and for $k = 400, 401, 402$.
  \item Passages through small decades agree with step-by-step computation, including the positions
        of all branch points.
  \item In base 3 the table has period $4$ and entry offsets $\{0, 2, 3, 6\}$, and it reproduces the
        transition table in the proof of Theorem~9.1 of~\cite{comma}.
\end{itemize}
